\documentclass[11pt]{article}

\usepackage[margin=1in]{geometry}
\usepackage{amsmath,amssymb,amsfonts}
\usepackage{array}
\usepackage{booktabs}
\usepackage{enumitem}
\usepackage{hyperref}
\usepackage{microtype}
\usepackage[T1]{fontenc}
\usepackage{lmodern}

\hypersetup{
  colorlinks=true,
  linkcolor=blue,
  urlcolor=blue,
  citecolor=blue,
}

\title{Epistemic Information Gain, Entropy, and Entropy Collapse in Reinforcement Learning}
\author{}
\date{}

\newcommand{\D}{\mathcal{D}}
\newcommand{\E}{\mathbb{E}}
\newcommand{\Env}{\mathcal{E}}
\newcommand{\A}{\mathcal{A}}
\newcommand{\Sref}{\mathcal{S}_{\mathrm{ref}}}
\newcommand{\KL}{\mathrm{KL}}
\newcommand{\Hent}{\mathrm{H}}
\newcommand{\I}{\mathrm{I}}

\begin{document}
\maketitle

\begin{abstract}
This note gives a first-principles account of how to measure information gained
by an RL agent after learning from a rollout on a single environment. The core
claim is simple: entropy reduction is meaningful only after naming the
uncertainty object whose entropy is being reduced. Entropy over actions measures
how decisive a policy is. Entropy over latent hypotheses measures epistemic
uncertainty. Expected epistemic information gain is expected posterior entropy
reduction, equivalently mutual information. In practical RL systems, this ideal
quantity is measured differently for a single non-Bayesian policy, an ensemble of
policies, and a Bayesian policy. The document ends with concrete estimators and
diagnostics for separating useful information gain from entropy collapse.
\end{abstract}

\section{The First Principle: Information Is Uncertainty Reduction}

Let $Z$ be an unknown quantity. It could be a hidden task identity, a reward
parameter, a transition parameter, a value-function hypothesis, or a policy
parameter. A belief over $Z$ is a probability distribution $p(z)$. The entropy of
that belief is
\begin{equation}
  \Hent[Z]
  =
  -\sum_z p(z)\log p(z),
\end{equation}
for a discrete variable. Entropy is high when the belief mass is spread across
many alternatives and low when the belief is concentrated.

If a new observation $Y$ is collected, the belief changes from $p(z)$ to
$p(z \mid y)$. The information gained from the observation is the reduction in
uncertainty:
\begin{equation}
  \Hent[Z] - \Hent[Z \mid Y=y].
\end{equation}
Before observing $Y$, the possible outcomes are unknown, so the expected
information gain is
\begin{equation}
  \I(Z;Y)
  =
  \Hent[Z] - \E_{y \sim p(y)}\Hent[Z \mid Y=y].
\end{equation}
This is mutual information. An equivalent expression is an expected KL
divergence:
\begin{equation}
  \I(Z;Y)
  =
  \E_{y \sim p(y)}
  \KL\left(p(z \mid y)\,\|\,p(z)\right).
\end{equation}

This equivalence is important in RL. The entropy-reduction view says
``uncertainty went down.'' The KL view says ``the observation changed the
belief.'' For continuous latent variables, differential entropy can be
coordinate-sensitive, so the KL form is often the cleaner measurement.

\section{Epistemic Information Gain From One Environment}

Let $\D$ be the agent's current data and let $e \in \Env$ be a particular
environment or level. A rollout from that environment is
\begin{equation}
  \tau = (s_0,a_0,r_0,s_1,\ldots,s_T).
\end{equation}
Let $\theta$ denote the latent object the agent is uncertain about. The exact
choice of $\theta$ depends on the scientific question:
\begin{itemize}[leftmargin=*]
  \item transition dynamics parameters;
  \item reward parameters;
  \item task identity;
  \item value-function hypotheses;
  \item policy parameters;
  \item the identity of which action is optimal in each state.
\end{itemize}

The expected epistemic information gain of rolling out on environment $e$ is
\begin{equation}
  \mathrm{EIG}(e)
  =
  \Hent[\theta \mid \D]
  -
  \E_{\tau \sim p(\tau \mid e,\D)}
  \Hent[\theta \mid \D,\tau].
\end{equation}
Equivalently,
\begin{equation}
  \mathrm{EIG}(e)
  =
  \I(\theta;\tau \mid \D,e)
  =
  \E_{\tau \sim p(\tau \mid e,\D)}
  \KL\left(p(\theta \mid \D,\tau)\,\|\,p(\theta \mid \D)\right).
\end{equation}

This is the clean target. It is epistemic because the entropy is over what is
true but unknown, not over stochastic actions or random outcomes. It is expected
because the rollout has not yet happened when choosing which environment to
sample.

\section{Nearby Methods for Scoring Informative Environments}

Several recent methods are close to the question in this note, but they place
the information signal in different objects: policy parameters, learned dynamics,
ensemble predictions, or task-policy relationships. This section summarizes four
representative methods and relates each to the ideal quantity
$\I(\theta;\tau \mid \D,e)$.

\subsection{PACE: Parameter Change as Realized Learning Progress}

PACE, proposed by Yuan et al., is the closest method to the single-policy
setting. It starts from the observation that many UED scores are indirect:
regret requires an optimal policy or an expert, value-error scores depend on the
critic, and Monte Carlo scores can have high variance. PACE instead scores a
level by the policy parameter change induced by training on that level.

Let $\theta^-$ be the policy parameters before a provisional update on level
$e$, and let $\theta_e^+$ be the parameters after that local update. The PACE
score is essentially
\begin{equation}
  S_{\mathrm{PACE}}(e)
  =
  \left\lVert \theta_e^+ - \theta^- \right\rVert_2^2.
\end{equation}
The motivation is a first-order Taylor approximation. If a local gradient update
on level $e$ is aligned with the level-specific objective, then the objective
improvement induced by that level is proportional to the squared norm of the
corresponding parameter update, up to a step-size constant. PACE therefore treats
parameter movement as an intrinsic, low-variance proxy for realized learning
progress.

Operationally, PACE collects a trajectory on a candidate level, performs a
temporary update to compute the induced parameter change, inserts high-scoring
levels into a replay buffer, and samples from that buffer with prioritization and
staleness correction. This maps naturally to the single non-Bayesian policy
regime in this note: it is not epistemic IG, because there is no posterior, but
it is a direct approximation to ``this environment caused the current learner to
move.''

\subsection{MaxInfoRL: Information Gain as an Intrinsic Reward}

MaxInfoRL, proposed by Sukhija et al., uses information gain inside online RL
exploration rather than as an external environment-selection score. The agent
learns a model of the unknown transition-and-reward function $f^\star$ and uses
an information-gain bonus to direct exploration toward transitions that reduce
uncertainty about $f^\star$.

For a transition query $(s,a)$, the ideal intrinsic reward is of the form
\begin{equation}
  r_{\mathrm{IG}}(s,a)
  =
  \I(\tilde{s}', f^\star \mid s,a,\D),
\end{equation}
where $\tilde{s}'$ denotes the predicted next transition target, possibly
including reward. In practice, the information gain does not generally have a
closed form, so MaxInfoRL uses an upper bound derived from epistemic uncertainty
in a learned model. With an ensemble or probabilistic model, the practical bonus
is high where model disagreement about the transition or reward is high.

MaxInfoRL then combines this intrinsic reward with maximum-entropy actor-critic
updates. In soft actor-critic notation, the policy is no longer pushed only
toward high extrinsic value and action entropy; it is also pushed toward states
and actions where the model expects high information gain. This method is a true
information-gain method with respect to dynamics/reward uncertainty, but its
unit of selection is a transition or behavior policy, not necessarily a reusable
environment level.

\subsection{MAX: Ensemble Disagreement for Active Exploration}

Model-Based Active Exploration (MAX), proposed by Shyam, Jaskowski, and Gomez,
also treats exploration as an active information-gathering problem. MAX trains
an ensemble of forward dynamics models and plans in a surrogate exploration MDP
whose reward is novelty of predicted transitions.

For a discrete next-state distribution, the ideal ensemble disagreement score is
\begin{equation}
  u_{\mathrm{MAX}}(s,a)
  =
  \Hent\left[
    \frac{1}{N}\sum_{i=1}^{N}p_i(s' \mid s,a)
  \right]
  -
  \frac{1}{N}\sum_{i=1}^{N}
  \Hent[p_i(s' \mid s,a)].
\end{equation}
This is the same entropy-of-the-mean minus mean-entropy structure used earlier
for action-space epistemic uncertainty. Here, however, it is applied to next
state predictions rather than action distributions. In continuous domains, MAX
replaces Shannon entropy with Renyi entropy and uses a Jensen-Renyi divergence
so that disagreement among Gaussian ensemble predictions can be computed
tractably.

MAX is therefore an ensemble approximation to epistemic information gain about
the dynamics. It is especially relevant to the ensemble-policy section of this
note because it shows the same mathematical decomposition in a different space:
not ``which action do policy hypotheses prefer?'' but ``which transition do
dynamics hypotheses predict?''

\subsection{ITTS: Information-Theoretic Task Selection for Meta-RL}

Information-Theoretic Task Selection (ITTS), proposed by Luna Gutierrez and
Leonetti, studies which training tasks should be selected for meta-RL. Unlike
PACE, MaxInfoRL, and MAX, ITTS is not primarily about online exploration within a
single environment. It asks which tasks are useful members of a meta-training
set.

ITTS keeps a candidate training set and a validation task set. It selects source
tasks that are both different from already selected tasks and relevant to the
validation tasks. Difference is measured by average KL divergence between the
policies learned for two candidate tasks, evaluated over states sampled from the
validation tasks:
\begin{equation}
  \delta(m_1,m_2)
  =
  \frac{1}{|\Sref|}
  \sum_{s \in \Sref}
  \KL\left(
    \pi_{m_1}^{\star}(\cdot \mid s)
    \,\|\,
    \pi_{m_2}^{\star}(\cdot \mid s)
  \right).
\end{equation}
Relevance is measured by adapting a policy from one task to another and checking
how the policy entropy changes over the validation distribution. A task is useful
when it is not redundant with selected tasks and transfers information to the
validation tasks.

For the present note, ITTS is important because it treats information as a
relationship between tasks mediated by policies. It is less directly applicable
to scoring a single rollout on one environment, but it suggests a useful
evaluation question: does learning on environment $e_i$ change the policy in a
way that transfers to a target task distribution, or only in a way that is
locally decisive on $e_i$?

\subsection{What Has Been Compared Carefully?}

The clearest metric-comparison work I found is No Regrets by Rutherford et al.
It asks whether common UED prioritization scores actually approximate the
quantities they are meant to approximate. In domains where regret can be
computed or probed more directly, they compare practical scores such as MaxMC
and positive value loss against regret, success rate, and learnability. Their key
finding is that practical regret approximations can correlate with success rate
rather than learnability, causing replay to spend experience on levels the agent
has already mastered.

PACE also compares against UED baselines such as DR, PLR, robust PLR, and ACCEL,
and explicitly discusses bias, variance, and computational cost of value-based,
Monte Carlo, regret, and marginal-benefit signals. However, it is primarily an
algorithmic comparison, not a full oracle study of every approximation against a
shared true information-gain target.

MaxInfoRL and MAX compare information-gain or disagreement-based exploration
against exploration baselines. ITTS compares selected task subsets against using
all or other subsets of meta-training tasks. I did not find a unified benchmark
that places PACE-style parameter change, MaxInfoRL-style model information gain,
MAX-style ensemble disagreement, and ITTS-style task-policy transfer under one
oracle information-gain evaluation.

\section{Entropy in RL}

RL contains several different entropies. Confusing them is the main source of
bad information-gain measurements.

\subsection{Policy Action Entropy}

For a stochastic policy $\pi(a \mid s)$, action entropy at state $s$ is
\begin{equation}
  \Hent_{\pi}(s)
  =
  -\sum_{a \in \A}\pi(a \mid s)\log \pi(a \mid s).
\end{equation}
This measures how spread out the policy's action distribution is. Low action
entropy means the policy is decisive. High action entropy means the policy is
diffuse or exploratory.

For a fixed probe state set $\Sref$, the average action entropy is
\begin{equation}
  \bar{\Hent}_{\pi}(\Sref)
  =
  \frac{1}{|\Sref|}
  \sum_{s \in \Sref}\Hent_{\pi}(s).
\end{equation}

\subsection{Trajectory Entropy}

A policy and environment together induce a distribution over trajectories:
\begin{equation}
  p(\tau \mid e,\pi).
\end{equation}
The entropy of this distribution measures how many possible trajectories the
agent-environment system can generate. It mixes policy stochasticity, transition
stochasticity, and state visitation diversity.

\subsection{Return Entropy}

The entropy of returns measures variability in cumulative reward:
\begin{equation}
  G(\tau)=\sum_{t=0}^{T}\gamma^t r_t.
\end{equation}
This is often about outcome uncertainty rather than epistemic uncertainty. A
fully understood environment can still have high return entropy if it is
intrinsically stochastic.

\subsection{Belief Entropy}

Belief entropy is entropy over latent hypotheses:
\begin{equation}
  \Hent[\theta \mid \D].
\end{equation}
This is the entropy that directly supports epistemic information gain. In a
Bayesian model it is explicit. In an ensemble it is approximated. In a single
non-Bayesian policy it is usually absent, so proxy metrics are required.

\section{Decisiveness Is Not Epistemic Information}

Suppose a policy has two actions at state $s$:
\begin{equation}
  \pi_{\mathrm{before}}(\cdot \mid s)=(0.5,0.5),
  \qquad
  \pi_{\mathrm{after}}(\cdot \mid s)=(0.95,0.05).
\end{equation}
The action entropy dropped. The agent became more decisive. But this does not
prove that the agent learned the truth about the environment. The same pattern
can arise from:
\begin{itemize}[leftmargin=*]
  \item correct learning from informative data;
  \item a changed entropy coefficient or temperature;
  \item optimizer instability;
  \item overfitting to a noisy rollout;
  \item a policy update that sharpened preferences without reducing model
  uncertainty.
\end{itemize}

The distinction is:
\begin{align}
  \text{Decisiveness}
  &=
  \text{concentration of }\pi(a \mid s),
  \\
  \text{Epistemic information}
  &=
  \text{reduction of uncertainty about }\theta.
\end{align}
Action entropy reduction is useful, but it is not automatically epistemic.

\section{Entropy Collapse in RL}

Entropy collapse is the rapid reduction of policy action entropy before the
agent has learned a robust, generalizable decision rule. It is not just
``becoming confident.'' It is becoming confident faster than the evidence
supports.

\subsection{Mechanisms}

Common mechanisms include:
\begin{itemize}[leftmargin=*]
  \item a small or decaying entropy bonus in policy-gradient training;
  \item large policy updates that overreact to high-variance advantages;
  \item sparse rewards that reinforce a lucky early trajectory;
  \item replay systems that repeatedly sample narrow states or levels;
  \item clipping or normalization choices that hide value-function uncertainty;
  \item deterministic action selection during data collection.
\end{itemize}

\subsection{Why Entropy Collapse Matters}

Entropy collapse can masquerade as information gain. If
\begin{equation}
  \bar{\Hent}_{\pi_{\mathrm{before}}}(\Sref)
  -
  \bar{\Hent}_{\pi_{\mathrm{after}}}(\Sref)
\end{equation}
is large, the policy became more concentrated. But that number alone cannot say
whether the rollout was informative. A bad rollout can make the agent confident
in the wrong action. In that case, entropy reduction is a warning sign, not a
success metric.

\subsection{Diagnostics}

Useful diagnostics for entropy collapse are:
\begin{itemize}[leftmargin=*]
  \item action entropy reduction on a fixed probe set;
  \item return or success-rate improvement on held-out environments;
  \item policy KL step size before and after learning;
  \item value-loss change on held-out trajectories;
  \item ensemble or posterior disagreement, when available;
  \item calibration of predicted values against realized returns.
\end{itemize}

A healthy update has a coherent pattern: the policy changes, held-out loss or
return improves, and epistemic disagreement decreases. An unhealthy collapse has
policy sharpening without corresponding held-out improvement or disagreement
reduction.

\section{Measuring After Learning on a Single Environment}

The experimental unit is one environment $e$. The measurement protocol should
compare the agent before and after learning from rollouts on $e$ while holding
the evaluation states fixed.

\subsection{Protocol}

\begin{enumerate}[leftmargin=*]
  \item Save the pre-learning agent state.
  \item Choose a fixed reference state set $\Sref$. This may be collected from a
  broad replay buffer, held-out levels, random starts, or a standardized
  evaluation distribution.
  \item Sample $K$ independent rollout batches
  $\tau_1,\ldots,\tau_K \sim p(\tau \mid e,\D)$.
  \item For each rollout batch $\tau_k$, copy the pre-learning agent and run the
  same learning procedure on $\D \cup \tau_k$.
  \item Measure before and after quantities on the same $\Sref$.
  \item Average the after-learning quantities over $k$.
\end{enumerate}

The fixed $\Sref$ matters. If entropy is measured only on states visited by each
rollout, the metric confounds information gain with a changed state
distribution.

\section{Single Non-Bayesian Policy}

A single non-Bayesian policy has parameters $\theta$ but no posterior
$p(\theta \mid \D)$. Therefore direct epistemic information gain is not
available. The practical target becomes learning progress: did data from
environment $e$ cause a useful, generalizing update?

\subsection{Action Entropy Reduction}

For rollout batch $\tau_k$, let $\pi_k^+$ be the policy after learning on
$\tau_k$ and $\pi^-$ the policy before learning. Define
\begin{equation}
  \Delta \Hent_{\mathrm{action}}(e)
  =
  \bar{\Hent}_{\pi^-}(\Sref)
  -
  \frac{1}{K}\sum_{k=1}^{K}
  \bar{\Hent}_{\pi_k^+}(\Sref).
\end{equation}
This measures decisiveness gain. It is not epistemic information gain by
itself.

\subsection{Policy Update Magnitude}

A direct measure of how much the rollout changed the policy is
\begin{equation}
  \Delta_{\mathrm{KL}}(e)
  =
  \frac{1}{K}\sum_{k=1}^{K}
  \frac{1}{|\Sref|}\sum_{s \in \Sref}
  \KL\left(\pi_k^+(\cdot \mid s)\,\|\,\pi^-(\cdot \mid s)\right).
\end{equation}
A symmetric version may be used when directional KL is unstable:
\begin{equation}
  \Delta_{\mathrm{SKL}}(e)
  =
  \frac{1}{2}
  \KL(\pi_k^+\,\|\,\pi^-)
  +
  \frac{1}{2}
  \KL(\pi^-\,\|\,\pi_k^+).
\end{equation}

\subsection{Probe Loss Reduction}

Let $\D_{\mathrm{probe}}$ be a fixed held-out dataset of states, actions,
returns, advantages, or value targets. Let $L(\theta;\D_{\mathrm{probe}})$ be a
chosen loss, such as PPO surrogate loss, behavior-cloning loss, TD error, value
loss, or return-prediction loss. Define
\begin{equation}
  \Delta L_{\mathrm{probe}}(e)
  =
  L(\theta^-;\D_{\mathrm{probe}})
  -
  \frac{1}{K}\sum_{k=1}^{K}
  L(\theta_k^+;\D_{\mathrm{probe}}).
\end{equation}
This is one of the best single-policy proxies because it asks whether learning
from $e$ improved performance on data not identical to the update rollout.

\subsection{Rollout Surprise}

If the agent has value predictions, define TD error
\begin{equation}
  \delta_t
  =
  r_t + \gamma V_{\theta^-}(s_{t+1}) - V_{\theta^-}(s_t).
\end{equation}
Then
\begin{equation}
  S_{\mathrm{TD}}(e)
  =
  \frac{1}{K}\sum_{k=1}^{K}
  \sum_{t \in \tau_k}|\delta_t|
\end{equation}
measures how surprising the rollout was under the current value function. High
surprise can identify candidate informative environments, but surprise without
subsequent probe improvement may be noise rather than information.

\subsection{Recommended Single-Policy Panel}

For a single non-Bayesian policy, report:
\begin{equation}
  \left(
  \Delta \Hent_{\mathrm{action}},
  \Delta_{\mathrm{KL}},
  \Delta L_{\mathrm{probe}},
  S_{\mathrm{TD}},
  \Delta \mathrm{return}_{\mathrm{heldout}}
  \right).
\end{equation}
Interpretation should be conservative:
\begin{itemize}[leftmargin=*]
  \item entropy reduction alone means decisiveness;
  \item policy KL means the rollout changed the policy;
  \item probe loss reduction means the update was useful;
  \item TD error means the rollout contradicted the critic;
  \item held-out return improvement is the final behavioral check.
\end{itemize}

\section{Ensemble Policy}

With $N$ policies, each policy can be treated as a hypothesis:
\begin{equation}
  \pi_1,\ldots,\pi_N.
\end{equation}
These may come from different random seeds, bootstrap resampling, different
checkpoints, or different training perturbations. This is not automatically
Bayesian, but it gives an empirical approximation to epistemic uncertainty.

\subsection{Predictive Action Distribution}

For each state $s$, define the ensemble-average action distribution:
\begin{equation}
  \bar{\pi}(a \mid s)
  =
  \frac{1}{N}\sum_{i=1}^{N}\pi_i(a \mid s).
\end{equation}
The total predictive action entropy is
\begin{equation}
  \Hent[\bar{\pi}(\cdot \mid s)].
\end{equation}
The average individual policy entropy is
\begin{equation}
  \frac{1}{N}\sum_{i=1}^{N}\Hent[\pi_i(\cdot \mid s)].
\end{equation}

\subsection{Disagreement as Epistemic Uncertainty}

The ensemble disagreement term is
\begin{equation}
  U_{\mathrm{epi}}(s)
  =
  \Hent[\bar{\pi}(\cdot \mid s)]
  -
  \frac{1}{N}\sum_{i=1}^{N}\Hent[\pi_i(\cdot \mid s)].
\end{equation}
This is high when individual policies are confident but disagree. It is low
when the policies agree, even if each policy is stochastic.

Average over the probe states:
\begin{equation}
  U_{\mathrm{epi}}(\Sref)
  =
  \frac{1}{|\Sref|}\sum_{s \in \Sref}U_{\mathrm{epi}}(s).
\end{equation}

\subsection{Information-Gain Proxy After Learning on One Environment}

Let the ensemble before learning be $\{\pi_i^-\}_{i=1}^{N}$. For rollout batch
$\tau_k$, train or update the ensemble to obtain $\{\pi_{i,k}^+\}_{i=1}^{N}$.
The ensemble epistemic-gain proxy is
\begin{equation}
  \widehat{\mathrm{EIG}}_{\mathrm{ens}}(e)
  =
  U_{\mathrm{epi}}^-(\Sref)
  -
  \frac{1}{K}\sum_{k=1}^{K}
  U_{\mathrm{epi},k}^+(\Sref).
\end{equation}

Separately, report decisiveness gain:
\begin{equation}
  \Delta \Hent_{\mathrm{individual}}(e)
  =
  \frac{1}{|\Sref|}\sum_{s \in \Sref}
  \left[
  \frac{1}{N}\sum_i \Hent[\pi_i^-(\cdot \mid s)]
  -
  \frac{1}{K}\sum_k
  \frac{1}{N}\sum_i \Hent[\pi_{i,k}^+(\cdot \mid s)]
  \right].
\end{equation}

The separation is essential:
\begin{itemize}[leftmargin=*]
  \item drop in individual entropy means the policies became sharper;
  \item drop in disagreement means the policies became more aligned;
  \item only the second is an epistemic proxy.
\end{itemize}

\section{Bayesian Policy}

A Bayesian policy maintains a posterior over parameters or hypotheses:
\begin{equation}
  p(\theta \mid \D).
\end{equation}
There may be one deployed agent, but its prediction is posterior predictive:
\begin{equation}
  p(a \mid s,\D)
  =
  \E_{\theta \sim p(\theta \mid \D)}
  \pi_{\theta}(a \mid s).
\end{equation}

\subsection{Action-Space Epistemic Uncertainty}

At a probe state $s$, total predictive action entropy is
\begin{equation}
  \Hent\left[
    \E_{\theta}\pi_{\theta}(\cdot \mid s)
  \right].
\end{equation}
The expected policy stochasticity is
\begin{equation}
  \E_{\theta}
  \Hent[\pi_{\theta}(\cdot \mid s)].
\end{equation}
The epistemic action uncertainty is the difference:
\begin{equation}
  U_{\mathrm{Bayes}}(s)
  =
  \Hent\left[
    \E_{\theta}\pi_{\theta}(\cdot \mid s)
  \right]
  -
  \E_{\theta}
  \Hent[\pi_{\theta}(\cdot \mid s)].
\end{equation}
This is the mutual information between the action and posterior hypothesis:
\begin{equation}
  U_{\mathrm{Bayes}}(s)
  =
  \I(a;\theta \mid s,\D).
\end{equation}

\subsection{Posterior Update From One Environment}

After rollout $\tau_k$ on environment $e$, the posterior becomes
\begin{equation}
  p(\theta \mid \D,\tau_k).
\end{equation}
The action-space epistemic information-gain estimator is
\begin{equation}
  \widehat{\mathrm{EIG}}_{\mathrm{Bayes-action}}(e)
  =
  U_{\mathrm{Bayes}}^-(\Sref)
  -
  \frac{1}{K}\sum_{k=1}^{K}
  U_{\mathrm{Bayes},k}^+(\Sref).
\end{equation}

This is an action-space proxy for epistemic information gain. The more direct
Bayesian quantity is the posterior KL:
\begin{equation}
  \widehat{\mathrm{EIG}}_{\mathrm{Bayes}}(e)
  =
  \frac{1}{K}\sum_{k=1}^{K}
  \KL\left(
    p(\theta \mid \D,\tau_k)
    \,\|\,
    p(\theta \mid \D)
  \right).
\end{equation}
When the posterior is represented by samples
$\theta_1,\ldots,\theta_N$, the computation resembles the ensemble estimator.
The interpretation is stronger: the samples approximate a posterior, not merely
a collection of independently trained policies.

\section{Expected Entropy Reduction Over Action Space}

If the desired metric is explicitly expected entropy reduction over actions,
the estimator is
\begin{equation}
  \mathrm{EER}_{\mathrm{action}}(e)
  =
  \bar{\Hent}_{\pi^-}(\Sref)
  -
  \frac{1}{K}\sum_{k=1}^{K}
  \bar{\Hent}_{\pi_k^+}(\Sref).
\end{equation}
This is valid as a measure of action-space entropy reduction. It should be
named as such. It becomes an epistemic metric only when the action distribution
is posterior predictive and the epistemic component is separated from expected
policy stochasticity.

\section{Practical Measurement Table}

\begin{center}
\small
\begin{tabular}{
  >{\raggedright\arraybackslash}p{0.17\linewidth}
  >{\raggedright\arraybackslash}p{0.34\linewidth}
  >{\raggedright\arraybackslash}p{0.37\linewidth}
}
\toprule
Setting & Main estimator & Interpretation \\
\midrule
Single-policy, non-Bayesian
&
$\Delta L_{\mathrm{probe}}$, $\Delta_{\mathrm{KL}}$,
$\Delta \Hent_{\mathrm{action}}$, TD error
&
Learning-progress proxy. Entropy reduction means decisiveness, not direct
epistemic information gain. \\
\addlinespace
Ensemble policy
&
$U_{\mathrm{epi}} =
\Hent[\bar{\pi}] - \frac{1}{N}\sum_i\Hent[\pi_i]$
&
Disagreement reduction approximates epistemic uncertainty reduction. \\
\addlinespace
Bayesian policy
&
$\KL(p(\theta \mid \D,\tau)\,\|\,p(\theta \mid \D))$ or
$\I(a;\theta \mid s,\D)$ reduction
&
Direct epistemic information gain when the posterior is well specified. \\
\bottomrule
\end{tabular}
\end{center}

\section{Recommended Reporting Template}

For each candidate environment $e$, report the following after $K$ rollout
batches and matched learning updates:
\begin{enumerate}[leftmargin=*]
  \item action entropy before and after on $\Sref$;
  \item policy KL before versus after on $\Sref$;
  \item probe loss before versus after on a fixed held-out dataset;
  \item held-out return or success-rate delta;
  \item ensemble or posterior disagreement reduction, if available;
  \item posterior KL, if a Bayesian posterior exists.
\end{enumerate}

Then classify the result:
\begin{itemize}[leftmargin=*]
  \item \textbf{Useful information gain}: probe loss improves, held-out behavior
  improves, and disagreement or posterior entropy decreases.
  \item \textbf{Decisiveness without evidence}: action entropy drops but probe
  loss and held-out performance do not improve.
  \item \textbf{Surprise without learning}: TD error or rollout loss is high,
  but subsequent probe improvement is absent.
  \item \textbf{Instability}: policy KL is large, entropy collapses, and held-out
  metrics worsen.
\end{itemize}

\section{Oracle Evaluation of Approximate Metrics}

To compare approximations carefully, define an expensive oracle that estimates
the actual utility of training on each candidate environment. The oracle need
not be deployable during normal training. Its purpose is to create a reference
ranking against which cheap metrics can be judged.

\subsection{Notation for the Oracle Experiment}

The oracle experiment has three layers: candidate environments, repeated trials
on each candidate, and a fixed target evaluation distribution. The notation is:

\begin{center}
\small
\begin{tabular}{
  >{\raggedright\arraybackslash}p{0.18\linewidth}
  >{\raggedright\arraybackslash}p{0.72\linewidth}
}
\toprule
Symbol & Meaning \\
\midrule
$\mathcal{E}_{\mathrm{cand}}$ & The pool of candidate environments being
ranked now. In a maze setting, these are the candidate levels available for
replay or training. \\
$M$ & Number of candidate environments in the pool. \\
$e_j$ & The $j$th candidate environment, with $j \in \{1,\ldots,M\}$. \\
$K$ & Number of independent oracle trials per candidate environment. Larger
$K$ gives a lower-variance oracle estimate. \\
$k$ & Trial index, with $k \in \{1,\ldots,K\}$. \\
$\theta^-$ & The same saved pre-learning agent parameters used at the start of
every oracle trial. \\
$\tau_{j,k}$ & The rollout batch collected from candidate environment $e_j$ in
trial $k$, using the pre-learning agent. \\
$U(\theta^-,\tau_{j,k})$ & The fixed learning procedure applied to the cloned
agent using rollout batch $\tau_{j,k}$. \\
$\theta_{j,k}^+$ & The post-learning parameters after training clone $k$ on
candidate environment $e_j$. \\
$p_{\mathrm{target}}(e)$ & The target evaluation distribution. This should be
held fixed and should represent the environments where generalization matters. \\
$J_{\mathrm{target}}(\theta)$ & Evaluation performance of agent $\theta$ on the
target distribution, such as mean return or success rate. \\
$O(e_j)$ & The expensive oracle score for candidate environment $e_j$. \\
$A(e_j)$ & A cheap approximation score for candidate environment $e_j$. \\
\bottomrule
\end{tabular}
\end{center}

The key design principle is that every candidate starts from the same
$\theta^-$. Otherwise, a candidate may look good simply because it was evaluated
from a better initial checkpoint. The target distribution must also stay fixed.
Otherwise, the score mixes ``this environment taught the agent'' with ``the
evaluation distribution changed.''

\subsection{Oracle Target}

Let $\mathcal{E}_{\mathrm{cand}}=\{e_1,\ldots,e_M\}$ be candidate environments.
For each $e_j$, start from the same pre-learning agent state $\theta^-$. Run
$K$ independent rollout-and-update trials:
\begin{equation}
  \theta_{j,k}^+
  =
  U(\theta^-, \tau_{j,k}),
  \qquad
  \tau_{j,k} \sim p(\tau \mid e_j,\theta^-).
\end{equation}
Then evaluate each updated agent on a fixed target distribution
$p_{\mathrm{target}}(e)$ or a fixed held-out suite:
\begin{equation}
  O_{\mathrm{return}}(e_j)
  =
  \frac{1}{K}\sum_{k=1}^{K}
  \left[
    J_{\mathrm{target}}(\theta_{j,k}^+)
    -
    J_{\mathrm{target}}(\theta^-)
  \right].
\end{equation}
This is a behavioral oracle: it measures actual downstream improvement caused by
training on $e_j$.

If a calibrated Bayesian posterior is available, one can also define a posterior
oracle:
\begin{equation}
  O_{\mathrm{Bayes}}(e_j)
  =
  \frac{1}{K}\sum_{k=1}^{K}
  \KL\left(
    p(\theta \mid \D,\tau_{j,k})
    \,\|\,
    p(\theta \mid \D)
  \right),
\end{equation}
or the corresponding reduction in posterior-predictive action disagreement on
$\Sref$. This oracle is closer to epistemic information gain, but it depends on
the posterior representation being meaningful.

\subsection{Approximation Metrics to Compare}

For each candidate environment $e_j$, compute cheap scores before committing to
expensive training:
\begin{align}
  A_{\mathrm{entropy}}(e_j)
  &=
  \Delta \Hent_{\mathrm{action}}(e_j),
  \\
  A_{\mathrm{policyKL}}(e_j)
  &=
  \Delta_{\mathrm{KL}}(e_j),
  \\
  A_{\mathrm{probe}}(e_j)
  &=
  \Delta L_{\mathrm{probe}}(e_j),
  \\
  A_{\mathrm{PACE}}(e_j)
  &=
  \left\lVert \theta_j^+ - \theta^- \right\rVert_2^2,
  \\
  A_{\mathrm{ens}}(e_j)
  &=
  U_{\mathrm{epi}}^-(\Sref)
  -
  U_{\mathrm{epi},j}^+(\Sref),
  \\
  A_{\mathrm{modelIG}}(e_j)
  &=
  \sum_{(s,a)\in \tau_j}
  \I(\tilde{s}', f^\star \mid s,a,\D).
\end{align}
These correspond, respectively, to decisiveness, behavioral movement,
held-out learning progress, PACE-style parameter change, ensemble epistemic
disagreement reduction, and MaxInfoRL/MAX-style model information.

\subsection{Worked Toy Example}

Consider three candidate maze levels:
\begin{align*}
  e_1 &: \text{an easy maze the agent mostly solves already}, \\
  e_2 &: \text{a medium maze with a new useful trap pattern}, \\
  e_3 &: \text{a very hard maze that is mostly unsolvable for the current agent}.
\end{align*}
Suppose the pre-learning agent has target-suite success
$J_{\mathrm{target}}(\theta^-)=50\%$. The oracle clones the same agent three
times per candidate, trains each clone on that candidate, and then evaluates on
the same held-out target maze suite. For simplicity, write the averaged result
directly:

\begin{center}
\small
\begin{tabular}{
  >{\raggedright\arraybackslash}p{0.12\linewidth}
  >{\raggedright\arraybackslash}p{0.25\linewidth}
  >{\raggedright\arraybackslash}p{0.28\linewidth}
  >{\raggedright\arraybackslash}p{0.22\linewidth}
}
\toprule
Level & After-training target success & Oracle value
$O_{\mathrm{return}}(e_j)$ & Interpretation \\
\midrule
$e_1$ & $52\%$ & $+2$ percentage points & Slightly useful. \\
$e_2$ & $65\%$ & $+15$ percentage points & Most useful. \\
$e_3$ & $48\%$ & $-2$ percentage points & Harmful or too hard. \\
\bottomrule
\end{tabular}
\end{center}

The oracle ranking is therefore
\begin{equation}
  e_2 \succ e_1 \succ e_3.
\end{equation}
Now compare cheap scores computed before paying for the full oracle:

\begin{center}
\footnotesize
\setlength{\tabcolsep}{3pt}
\begin{tabular}{@{}
  >{\raggedright\arraybackslash}p{0.22\linewidth}
  >{\raggedright\arraybackslash}p{0.14\linewidth}
  >{\raggedright\arraybackslash}p{0.14\linewidth}
  >{\raggedright\arraybackslash}p{0.14\linewidth}
  >{\raggedright\arraybackslash}p{0.18\linewidth}
@{}}
\toprule
Approximation & $A(e_1)$ & $A(e_2)$ & $A(e_3)$ & Proxy pick \\
\midrule
Action entropy reduction & $0.10$ & $0.25$ & $0.60$ & $e_3$ \\
TD error magnitude & $0.20$ & $0.45$ & $0.90$ & $e_3$ \\
Probe loss reduction & $0.03$ & $0.18$ & $-0.04$ & $e_2$ \\
PACE parameter change & $0.05$ & $0.30$ & $0.12$ & $e_2$ \\
Ensemble disagreement reduction & $0.02$ & $0.22$ & $0.01$ & $e_2$ \\
\bottomrule
\end{tabular}
\end{center}

In this toy example, entropy reduction and TD error pick $e_3$ because the hard
maze causes surprise and sharp policy changes. But the oracle shows that
training on $e_3$ hurts target performance. Probe loss reduction, PACE parameter
change, and ensemble disagreement reduction pick $e_2$, matching the oracle
best environment. The conclusion would be that entropy reduction and TD error
are detecting disruption or difficulty, while the latter three proxies better
track useful learning for this checkpoint and environment family.

\subsection{Comparison Criteria}

The oracle comparison should not only ask which score gives the best final
policy. It should ask whether the approximation tracks the oracle at the level
of candidate environments:
\begin{itemize}[leftmargin=*]
  \item \textbf{Rank correlation}: Spearman or Kendall correlation between
  $A(e_j)$ and $O(e_j)$ across candidate environments.
  \item \textbf{Top-$q$ overlap}: fraction of oracle top environments recovered
  by the approximation.
  \item \textbf{Selection regret}: oracle value lost by choosing the top
  approximation-ranked environment instead of the top oracle-ranked
  environment.
  \item \textbf{Calibration}: whether score magnitude predicts expected
  downstream improvement, not just rank.
  \item \textbf{Robustness}: whether the correlation survives different seeds,
  training phases, environment families, and target evaluation distributions.
\end{itemize}

This design turns the conceptual question into a falsifiable experiment. If
action entropy reduction ranks environments poorly against
$O_{\mathrm{return}}$, it is merely a decisiveness metric. If PACE or probe-loss
reduction tracks $O_{\mathrm{return}}$ well, it is a useful single-policy proxy.
If ensemble disagreement reduction tracks $O_{\mathrm{Bayes}}$ or
$O_{\mathrm{return}}$, it is a stronger epistemic proxy. If model information
gain tracks only novelty but not downstream improvement, it may be valuable for
exploration but insufficient as a curriculum score.

\section{Conclusion}

The first-principles definition of epistemic information gain is expected
posterior entropy reduction over the latent quantity the agent is uncertain
about. In RL, action entropy is tempting because it is easy to compute, but it
mostly measures decisiveness. Entropy collapse is the failure mode where this
decisiveness appears before reliable knowledge.

For a single non-Bayesian policy, epistemic information gain must be approximated
by learning-progress metrics such as probe loss reduction, policy-update KL, TD
error, and held-out behavioral improvement. For an ensemble, disagreement
reduction gives a useful epistemic proxy. For a Bayesian policy, posterior KL or
posterior-predictive mutual information gives the cleanest measurement. Recent
work supplies complementary approximations: PACE uses parameter change, MaxInfoRL
uses model information gain, MAX uses ensemble disagreement over dynamics, and
ITTS uses information-theoretic policy relationships between tasks. The strongest
way to compare these approximations is to construct an expensive oracle that
measures actual downstream improvement or true posterior change, then evaluate
which cheap score best predicts the oracle ranking.

\section*{References}

\begin{itemize}[leftmargin=*]
  \item Shannon, C. E. (1948). A mathematical theory of communication.
  \item Lindley, D. V. (1956). On a measure of the information provided by an
  experiment.
  \item MacKay, D. J. C. (1992). Information-based objective functions for
  active data selection.
  \item Gal, Y., Islam, R., and Ghahramani, Z. (2017). Deep Bayesian active
  learning with image data.
  \item Haarnoja, T., Zhou, A., Abbeel, P., and Levine, S. (2018). Soft
  actor-critic: Off-policy maximum entropy deep reinforcement learning with a
  stochastic actor.
  \item Schulman, J., Wolski, F., Dhariwal, P., Radford, A., and Klimov, O.
  (2017). Proximal policy optimization algorithms.
  \item Yuan, F., Yin, Q., Shen, S., Xie, Y., Yang, J., Qin, L., Zeng, J., and
  Li, Q. (2026). PACE: Parameter Change for Unsupervised Environment Design.
  \href{https://arxiv.org/abs/2605.01358}{arXiv:2605.01358}.
  \item Sukhija, B., Coros, S., Krause, A., Abbeel, P., and Sferrazza, C.
  (2025). MaxInfoRL: Boosting exploration in reinforcement learning through
  information gain maximization. ICLR.
  \href{https://sukhijab.github.io/projects/maxinforl/}{Project page}.
  \item Shyam, P., Jaskowski, W., and Gomez, F. (2019). Model-Based Active
  Exploration. ICML.
  \href{https://proceedings.mlr.press/v97/shyam19a.html}{PMLR 97}.
  \item Luna Gutierrez, R., and Leonetti, M. (2020). Information-theoretic Task
  Selection for Meta-Reinforcement Learning. NeurIPS.
  \href{https://arxiv.org/abs/2011.01054}{arXiv:2011.01054}.
  \item Rutherford, A., Beukman, M., Willi, T., Lacerda, B., Hawes, N., and
  Foerster, J. (2024). No Regrets: Investigating and Improving Regret
  Approximations for Curriculum Discovery. NeurIPS.
  \href{https://arxiv.org/abs/2408.15099}{arXiv:2408.15099}.
\end{itemize}

\end{document}
